\scanpage{136}\mspage{10}
% archive sheet number: 135
par $\Delta(f)$ dans $\unsure{{}_{X\backslash}M|{}_X M}$), le morphisme $X\xrightarrow{i_1(f)}X_1(f)$ est le morphisme \unsure{costructural}, et le morphisme $X_1(f)\xrightarrow{p_1(f)}Y$ étant \unsure{défini} \deleted{\unsure{qui}} par \unsure{coproduit|composé} des $p_D$. En somme, on fait \unsure{au} plus \unsure{économique} (et \unsure{au} plus naturel) pour assurer la validité de 1.17, sous la forme équivalente (1.15.). D'autre part, par définition \inserted{de $C'_0$, $C''_0$}, les $i_D$ sont dans $C'_0$, et $p_1(f)$ est dans $C''_0$, ce qui est 1.20.
\margin{$\Vert$ Il faudrait quand même expliciter la dépendance fonctorielle de $X_1(f)$ \unsure{p.r.|p.v.} à $f$ !!!}

Pour résumer :

\uline{Théorème 1} Soit $M$ une $\unsure{\mathcal U|U}$-catégorie stable par $\mathcal U$-limites inductives filtrantes, $C_0\subset\mathrm{Fl}\,M$ un petit ensemble de flèches de $M$. On suppose les flèches $i\in C_0$ \unsure{coquarrables}. Cela assure la \deleted{\unsure{exist}} définition d'un \inserted{endo}foncteur $f\mapsto p_1(f)$ de $\uline{\mathrm{Fl}}(M)$ \deleted{\unsure{sur}}, et d'un hom. fonctoriel \inserted{$i_1:$} $\mathrm{id}_{\uline{\mathrm{Fl}}(M)}\longrightarrow\unsure{P_1|p_1}$, i.e. $i_1(f):\deleted{\unsure{\to}}f\longrightarrow p_1(f)$, induisant l'identité sur les objets buts de $f$ et de $p_1(f)$,
\diagram{\begin{tikzcd}[ampersand replacement=\&,column sep=small,row sep=large]
X\arrow[rr,"i_1(f)"]\arrow[dr,"f"']\&\&X_1(f)\arrow[dl,"p_1(f)"]\\
\&Y\&
\end{tikzcd}}
d'où par itération un système inductif ordinal d'endofoncteurs $f\mapsto p_\alpha(f)$ de $\mathrm{Fl}(M)$, et \deleted{\unsure{d'hom. fonct.}} d'homomorphismes $i_\alpha:\mathrm{id}_{\mathrm{Fl}(M)}\longrightarrow P_\alpha$ i.e.
\diagram{\begin{tikzcd}[ampersand replacement=\&,column sep=small,row sep=large]
X\arrow[rr,dash,"i_\alpha(f)"]\arrow[dr,"f"']\&\&X_\alpha(f)\arrow[dl,"p_\alpha(f)"]\\
\&Y\&
\end{tikzcd}}
\unsure{formant|forment} un système \unsure{inductif} ordinal de factorisations \unsure{des|de} $f$, fonctorielles en $f$. \deleted{Ceci posé \unsure{Il existe} \unsure{soit de}}
